% sec:adm-equations — The ADM Evolution and Constraint Equations
\section{The ADM Evolution and Constraint Equations}
\label{sec:adm-equations}

The spatial metric $\spatialmetric$ and the extrinsic curvature $\extrinsic$
carry the gravitational initial data, but Einstein's equations have not yet
told us which combinations of these fields are freely specifiable and which
are constrained, nor how the unconstrained data evolve from one slice to the
next.  The answer comes from decomposing the spacetime Riemann tensor into
parts tangent and normal to $\Sigma_t$ --- the Gauss and Codazzi equations
of hypersurface geometry --- and substituting into Einstein's field
equations.  The result is a constrained evolution system: four constraint
equations that the initial data must satisfy on every slice, and twelve
evolution equations that propagate $\spatialmetric$ and $\extrinsic$
forward in time.  The system is the gravitational analogue of the
Maxwell constraint--evolution split, where $\nabla \cdot \mathbf{E} = \rho$
constrains the initial electric field and Faraday's law evolves it ---
except that here both the ``field'' ($\spatialmetric$) and its ``velocity''
($\extrinsic$) are constrained.  The analogy breaks down in two ways: the
Maxwell constraints are linear in the fields, whereas the ADM constraints
are nonlinear; and in electrodynamics only the field is constrained
($\nabla \cdot \mathbf{E} = \rho$), whereas here $\spatialmetric$ and
$\extrinsic$ are jointly constrained by coupled elliptic equations.
\physnote{In electrodynamics, the divergence constraints propagate
automatically if satisfied initially.  The same holds for the ADM
constraints --- but only at the continuum level.  Numerical violations of
the constraints can grow, motivating the reformulations of
\cref{ch:bssn-ccz4}.}

\paragraph{The Gauss equation.}

The Gauss equation relates the intrinsic curvature of $\Sigma_t$ to the
spacetime curvature projected entirely onto the slice.  Projecting all four
indices of the spacetime Riemann tensor onto $\Sigma_t$ with the spatial
projector $\gamma^\mu{}_\nu$ (\cref{eq:adm-projector}) gives
%
\begin{equation}
  \gamma^\alpha{}_\mu\,\gamma^\beta{}_\nu\,
  \gamma^\gamma{}_\rho\,\gamma^\delta{}_\sigma\,
  R_{\alpha\beta\gamma\delta}
  = {}^{(3)}\!R_{\mu\nu\rho\sigma}
    + K_{\mu\rho}\,K_{\nu\sigma}
    - K_{\mu\sigma}\,K_{\nu\rho} \,,
  \label{eq:adm-gauss}
\end{equation}
%
where ${}^{(3)}\!R_{\mu\nu\rho\sigma}$ is the Riemann tensor built from
$\spatialmetric$ and its spatial Christoffel symbols
$\spatialchristoffel{i}{j}{k}$ alone.  The two $K$-squared terms encode
the contribution of the embedding: even if the spatial geometry is
intrinsically flat (${}^{(3)}\!R_{\mu\nu\rho\sigma} = 0$), a curved
embedding ($\extrinsic \neq 0$) generates nonzero spacetime curvature
within the slice.  \Cref{eq:adm-gauss} is the higher-dimensional
generalisation of Gauss's \emph{Theorema Egregium} --- the remarkable
theorem that the Gaussian curvature of a surface can be computed from
measurements within the surface alone, without reference to the ambient
space.  Here the spatial Riemann tensor plays the role of the intrinsic
curvature, and the extrinsic curvature terms capture the corrections from
the embedding.
\histnote{Gauss proved his Theorema Egregium (``remarkable theorem'') in
1827 for surfaces in $\mathbb{R}^3$.  The generalisation to hypersurfaces
in pseudo-Riemannian manifolds is due to the work of Eisenhart and others
in the early twentieth century.}

\paragraph{The Codazzi equation.}

Projecting three indices of the Riemann tensor onto $\Sigma_t$ and one
along the normal gives
%
\begin{equation}
  \gamma^\alpha{}_\mu\,n^\beta\,
  \gamma^\gamma{}_\nu\,\gamma^\delta{}_\rho\,
  R_{\alpha\beta\gamma\delta}
  = D_\mu K_{\nu\rho} - D_\nu K_{\mu\rho} \,,
  \label{eq:adm-codazzi}
\end{equation}
%
where $D_i$ is the spatial covariant derivative introduced in
\cref{sec:adm-extrinsic-curvature}.  The Codazzi equation constrains how
the extrinsic curvature varies across the slice: you cannot freely
prescribe how the bending varies from point to point --- the variation must
be consistent with the spacetime curvature threading through the slice.
Physically, the Codazzi equation is the integrability condition ensuring
that the slice fits smoothly into the ambient spacetime.

\paragraph{The Ricci equation.}

The final projection --- two indices along the normal and two tangent to the
slice --- yields the Ricci equation (sometimes called the Mainardi equation):
%
\begin{equation}
  \gamma^\alpha{}_\mu\,n^\beta\,\gamma^\gamma{}_\nu\,n^\delta\,
  R_{\alpha\beta\gamma\delta}
  = -\lie{n} K_{\mu\nu} - \frac{1}{\lapse}\,D_\mu D_\nu \lapse
    + K_{\mu\lambda}\,K^\lambda{}_\nu \,.
  \label{eq:adm-ricci}
\end{equation}
%
This equation determines how the extrinsic curvature changes along the
normal direction.  The $D_\mu D_\nu \lapse / \lapse$ term reflects the
dependence on the gauge choice: different lapse functions produce different
rates of change of $\extrinsic$, even for the same physical spacetime.
When combined with the Einstein equation, this becomes the $\extrinsic$
evolution equation (\cref{eq:adm-kij-evolution}) that propagates the
extrinsic curvature forward in time.
\warnnote{The Ricci equation (\cref{eq:adm-ricci}) is distinct from the
Ricci \emph{tensor} $\ricci$.  The name refers to Gregorio Ricci-Curbastro
only indirectly; the equation is more commonly attributed to Mainardi and
Codazzi in the differential geometry literature.}

\paragraph{The Hamiltonian constraint.}

Contracting the Gauss equation (\cref{eq:adm-gauss}) on its first and
third indices with $\gamma^{\mu\rho}$ produces the spatial Ricci tensor
$\spatialricci$ on the left and a combination of $K$-squared terms on the
right.  A second contraction with $\inversespatialmetric$ gives a scalar
identity.  Substituting into the $\einstein\,n^\mu\,n^\nu$ projection of
Einstein's equation $\einstein = 8\pi\,\stressenergy$ yields the
\emph{Hamiltonian constraint}:
%
\begin{equation}
  \spatialricciscalar + \meanextrinsic^2
    - K_{ij}\,K^{ij} = 16\pi\,\rho \,,
  \label{eq:adm-hamiltonian}
\end{equation}
%
where $\spatialricciscalar = \inversespatialmetric\,\spatialricci$ is the
spatial Ricci scalar, $\meanextrinsic = \inversespatialmetric\,\extrinsic$
is the mean curvature (\cref{eq:adm-mean-curvature}), and $\rho$ is the
energy density measured by Eulerian observers:
$\rho = \stressenergy\,n^\mu\,n^\nu$.  For vacuum spacetimes the right-hand
side vanishes and the constraint becomes
$\spatialricciscalar + \meanextrinsic^2 - K_{ij}\,K^{ij} = 0$.  This is
an elliptic equation --- it contains no time derivatives, only spatial
derivatives of $\spatialmetric$ (through $\spatialricciscalar$) and
algebraic combinations of $\extrinsic$.  It must hold on every slice, not
just the initial one; any solution of the evolution equations that drifts
away from satisfying \cref{eq:adm-hamiltonian} has developed a constraint
violation.
\physnote{The Hamiltonian constraint is the relativistic generalisation of
the Newtonian Poisson equation $\laplacian\Phi = 4\pi\rho$
(\cref{sec:gr-efe}).  In the weak-field, slow-motion limit the
$K$-squared terms vanish and $\spatialricciscalar$ reduces to spatial
derivatives of the Newtonian potential, recovering Poisson's equation.}

\paragraph{The momentum constraint.}

The mixed projection $\einstein\,\gamma^{\mu}{}_i\,n^\nu = 8\pi\,S_i$,
where $S_i = -\gamma_{i\mu}\,T^{\mu\nu}\,n_\nu$ is the momentum density
measured by Eulerian observers, contracts the Codazzi equation
(\cref{eq:adm-codazzi}) to give
%
\begin{equation}
  D_j\bigl(K^{ij} - \inversespatialmetric\,\meanextrinsic\bigr) = 8\pi\,S^i \,.
  \label{eq:adm-momentum}
\end{equation}
%
Like the Hamiltonian constraint, this is a purely spatial equation with no
time derivatives.  It constrains the divergence of the extrinsic curvature:
the trace-corrected combination
$K^{ij} - \gamma^{ij}\,\meanextrinsic$ must have a specific divergence
set by the matter momentum density.  In vacuum the right-hand side
vanishes, and the extrinsic curvature must be divergence-free in this
trace-corrected sense.  The three components of \cref{eq:adm-momentum}
together with the single Hamiltonian constraint
(\cref{eq:adm-hamiltonian}) constitute four constraint equations --- the
same count as the four gauge degrees of freedom ($\lapse$, $\shift$)
removed from the twelve components of $\spatialmetric$ and $\extrinsic$.
\implnote{Monitoring the $L^2$ norms of the Hamiltonian and momentum
constraint violations during evolution provides the primary diagnostic for
numerical accuracy.  A well-resolved simulation shows constraint violations
converging to zero with increasing resolution at the expected order of the
spatial discretisation.}

\paragraph{The $\extrinsic$ evolution equation.}

The evolution equation for $\spatialmetric$ was already derived in
\cref{sec:adm-extrinsic-curvature} (\cref{eq:adm-metric-evolution}):
$\partial_t\gamma_{ij} = -2\lapse\,\extrinsic + D_i\beta_j + D_j\beta_i$.
The companion equation for $\extrinsic$ follows from the Ricci equation
(\cref{eq:adm-ricci}) and the spatial projection of Einstein's equation.
The full derivation combines the Ricci equation (\cref{eq:adm-ricci}) with
the spatial projection of the Einstein equation, converts the Lie
derivative along $n^\mu$ to a coordinate time derivative using
$\lie{n} = (1/\lapse)(\partial_t - \lie{\beta})$ as in
\cref{sec:adm-extrinsic-curvature}, and eliminates the spacetime Riemann
projection in favour of spatial curvature and matter terms.  The
calculation is the most involved in the chapter; we state the result and
focus on its structure:
%
\begin{equation}
  \begin{split}
  \partial_t \extrinsic
  &= -D_i D_j \lapse
     + \lapse\bigl(\spatialricci
       + \meanextrinsic\,\extrinsic - 2\,K_{ik}\,K^k{}_j\bigr) \\
  &\quad + \beta^k\,D_k\,\extrinsic
     + K_{ik}\,D_j\beta^k + K_{kj}\,D_i\beta^k \\
  &\quad + 4\pi\,\lapse\bigl[
       (\gamma^{kl}\,S_{kl} - \rho)\,\gamma_{ij} - 2\,S_{ij}\bigr] \,,
  \end{split}
  \label{eq:adm-kij-evolution}
\end{equation}
%
where $S_{ij} = \gamma_{i\mu}\,\gamma_{j\nu}\,T^{\mu\nu}$ is the spatial
stress tensor.  The first line contains the geometric content: the spatial
derivatives of the lapse (from the Ricci equation), the spatial Ricci
tensor (intrinsic curvature of the slice), and the nonlinear $K$-squared
terms (the embedding contribution).  The second line contains the Lie
derivative along the shift, accounting for the coordinate drift between
slices --- exactly the same structure as the shift terms in
\cref{eq:adm-metric-evolution}.  The third line couples gravity to matter.

In vacuum ($\stressenergy = 0$) the matter terms vanish, and the evolution
equation depends only on $\spatialmetric$, $\extrinsic$, $\lapse$, and
$\shift$.  The spatial Ricci tensor $\spatialricci$ involves second spatial
derivatives of $\spatialmetric$ through the three-dimensional Christoffel
symbols, making the system second-order in space and first-order in time.
\coderef{ADMSpacetime}{spatialChristoffel}

\paragraph{Structure of the ADM system.}

The complete ADM system consists of:
%
\begin{itemize}
  \item \textbf{Twelve evolution variables:} the six independent components
    of $\spatialmetric$ and the six of $\extrinsic$.
  \item \textbf{Twelve evolution equations:}
    \cref{eq:adm-metric-evolution,eq:adm-kij-evolution}.
  \item \textbf{Four constraints:} the Hamiltonian constraint
    (\cref{eq:adm-hamiltonian}) and the three components of the momentum
    constraint (\cref{eq:adm-momentum}).
  \item \textbf{Four gauge functions:} $\lapse$ and the three components of
    $\shift$, freely specifiable.
\end{itemize}
%
The constraints involve no time derivatives and must be satisfied on every
slice.  Solving twelve evolution equations for sixteen variables (plus four
gauge functions) seems under-determined --- why is this consistent?  A
fundamental property of the Einstein equations guarantees that if
the constraints are satisfied on the initial slice and the evolution
equations are solved exactly, then the constraints remain satisfied on all
subsequent slices --- the constraints \emph{propagate}.  This follows from
the contracted Bianchi identity $\cd{\mu}G^{\mu\nu} = 0$
(\cref{sec:dg-curvature}): the divergence-free nature of the Einstein
tensor implies that the time derivatives of the constraint quantities can
be expressed in terms of the constraints themselves, so vanishing
constraints remain vanishing.

\paragraph{Weak hyperbolicity.}

The ADM system as written suffers from a subtle but fatal deficiency for
numerical work: it is only \emph{weakly hyperbolic}.  A PDE system is
strongly hyperbolic if its principal symbol --- the matrix of highest-order
spatial derivative coefficients, evaluated along an arbitrary spatial
direction --- has a complete set of real eigenvalues and a complete basis of
eigenvectors.  Strong hyperbolicity guarantees well-posedness: solutions
depend continuously on initial data, and perturbations grow at most
exponentially in time.

The ADM evolution equation for $\extrinsic$ (\cref{eq:adm-kij-evolution})
contains the spatial Ricci tensor $\spatialricci$, which involves second
spatial derivatives of $\spatialmetric$ through the three-dimensional
Christoffel symbols.  Analysing the principal part of the combined system
(\cref{eq:adm-metric-evolution,eq:adm-kij-evolution}) reveals that the
principal symbol has the correct real eigenvalues (the physical
characteristic speeds $0$ and $\pm\lapse$), but does \emph{not} have a
complete set of eigenvectors.  The missing eigenvectors correspond to
gauge modes that propagate at zero speed.  As a result, the system is only
weakly hyperbolic: the initial-value problem is not well-posed in the
standard Hadamard sense, and numerical solutions can develop
resolution-dependent instabilities that grow polynomially --- not
exponentially --- with time.  In practice, this means that doubling the
resolution does not improve the solution but can actually make it worse by
resolving previously invisible unstable modes.
\warnnote{Weak hyperbolicity does not mean the physics is wrong --- the
constraint surface of the ADM system is perfectly well-defined, and the
physical degrees of freedom propagate correctly.  The problem is entirely
with the gauge sector: the coordinate degrees of freedom lack the
structure needed for well-posed numerical evolution.}

The resolution is to reformulate the equations so that all modes ---
physical and gauge --- propagate with a complete set of eigenvectors.  The
BSSN and CCZ4 formulations of \cref{ch:bssn-ccz4} achieve this by
introducing auxiliary variables (conformal connection functions
$\conformalconnection$) that promote the problematic zero-speed gauge modes
to well-behaved propagating modes, and by adding constraint-damping terms
that actively suppress numerical constraint violations.

\paragraph{Codebase note.}

The \texttt{ADMSpacetime} typeclass (\cref{sec:adm-lapse-shift}) provides
all the fields appearing in the ADM system: the lapse, shift, spatial
metric, inverse spatial metric, extrinsic curvature, spatial Christoffel
symbols, lapse gradient, and shift Jacobian.  For analytic spacetimes these
are closed-form expressions; for numerical evolutions they are grid
functions updated at each time step by the BSSN or CCZ4 right-hand sides.
The evolution code does not integrate the ADM equations directly ---
it uses the strongly hyperbolic reformulations of \cref{ch:bssn-ccz4} ---
but the ADM decomposition provides the conceptual and notational foundation
on which those reformulations build.
%
\begin{lstlisting}[language=Haskell]
  extrinsicCurvature :: s -> Point s -> Sym3 'Covariant
  spatialChristoffel :: s -> Point s -> Christoffel3
  lapseGradient      :: s -> Point s -> Vec3 'Covariant
  shiftJacobian      :: s -> Point s -> Mat3
\end{lstlisting}
%
The \texttt{lapseGradient} and \texttt{shiftJacobian} methods correspond to
$D_i\lapse$ and $\partial_j\beta^i$ appearing in the evolution equations
--- the spatial derivatives of the gauge variables that a numerical code
must compute at every grid point and every time step.
\xrefnote{The detailed analysis of the ADM principal symbol and its
incomplete eigenvector structure is carried out in
\cref{sec:bssn-adm-illposed}, where it motivates the BSSN conformal
decomposition.}